\documentclass{article}
\usepackage{amsmath,amssymb,amsthm}
\usepackage{algorithm,algorithmic}
\usepackage{enumitem}
\usepackage{hyperref}
\newtheorem{theorem}{Theorem}
\newtheorem{lemma}{Lemma}
\newtheorem{definition}{Definition}
\newtheorem{remark}{Remark}
\begin{document}

\section*{Proximal Gradient Method (PGM)}

\section*{Mathematical Formulation}
Let 
\[
F(x)\;:=\;f(x)+g(x),\qquad x\in\mathbb{R}^{d},
\]
where  

\begin{itemize}
  \item $f:\mathbb{R}^{d}\rightarrow\mathbb{R}$ is convex and $L$‑smooth, i.e.
    \[
    \|\nabla f(x)-\nabla f(y)\|\le L\|x-y\|\quad\forall x,y,
    \]
  \item $g:\mathbb{R}^{d}\rightarrow\mathbb{R}\cup\{+\infty\}$ is proper, closed and convex (possibly nonsmooth).
\end{itemize}
For a stepsize $\eta>0$ the proximal operator of $g$ is
\[
\operatorname{prox}_{\eta g}(v)\;:=\;\arg\min_{u\in\mathbb{R}^{d}}
\Big\{ g(u)+\frac{1}{2\eta}\|u-v\|^{2}\Big\}.
\]

The **proximal gradient iteration** reads
\begin{equation}
\boxed{
x_{k+1}= \operatorname{prox}_{\eta g}\!\bigl(x_{k}-\eta \nabla f(x_{k})\bigr)},
\qquad k=0,1,2,\dots
\label{eq:pgm}
\end{equation}
A stepsize satisfying $0<\eta\le \frac{1}{L}$ guarantees convergence.

\section*{Pseudocode}
\begin{algorithm}[H]
\caption{Proximal Gradient Method}
\begin{algorithmic}[1]
\REQUIRE initial point $x_{0}\in\mathbb{R}^{d}$, stepsize $\eta\in(0,1/L]$,
         maximum iterations $K$.
\FOR{$k=0,\dots,K-1$}
    \STATE $y_{k}\gets x_{k}-\eta\nabla f(x_{k})$ \hfill\Comment{gradient step}
    \STATE $x_{k+1}\gets \operatorname{prox}_{\eta g}(y_{k})$ \hfill\Comment{prox step}
\ENDFOR
\RETURN $x_{K}$
\end{algorithmic}
\end{algorithm}

\section*{Dry Run Example}
Consider the scalar problem
\[
f(w)=(w-3)^{2},\qquad g\equiv0,\qquad \eta=0.1,\qquad w_{0}=0.
\]
Since $g\equiv0$, $\operatorname{prox}_{\eta g}(v)=v$.

\begin{enumerate}[label=Iteration \arabic*:, leftmargin=*]
\item $k=0$  
  \[
  \nabla f(w_{0})=2(w_{0}-3)=-6.
  \]
  Gradient step: $y_{0}=w_{0}-\eta\nabla f(w_{0})=0-0.1(-6)=0.6$.  
  Prox step: $w_{1}=y_{0}=0.6$.  
  Objective: $F(w_{1})=f(0.6)=(0.6-3)^{2}=5.76$.
\item $k=1$  
  \[
  \nabla f(w_{1})=2(0.6-3)=-4.8.
  \]
  $y_{1}=0.6-0.1(-4.8)=1.08$, $w_{2}=1.08$, $F(w_{2})=(1.08-3)^{2}=3.6864$.
\item $k=2$  
  \[
  \nabla f(w_{2})=2(1.08-3)=-3.84.
  \]
  $y_{2}=1.08-0.1(-3.84)=1.464$, $w_{3}=1.464$, $F(w_{3})=(1.464-3)^{2}=2.3665$.
\end{enumerate}
The sequence $\{w_{k}\}$ moves toward the minimiser $w^{\star}=3$.

\newpage
\section*{Assumptions}
\begin{enumerate}[label=(A\arabic*)]
\item \textbf{Smoothness of $f$.} $f$ is convex and differentiable with $L$‑Lipschitz gradient:
\[
\|\nabla f(x)-\nabla f(y)\|\le L\|x-y\|\quad\forall x,y\in\mathbb{R}^{d}.
\tag{A1}
\]
\item \textbf{Convexity of $g$.} $g$ is proper, closed and convex.
\tag{A2}
\item \textbf{Stepsize.} The constant stepsize satisfies
\[
0<\eta\le\frac{1}{L}.
\tag{A3}
\]
\item \textbf{Existence of a minimiser.} There exists $x^{\star}\in\arg\min_{x}F(x)$ and the optimal value $F^{\star}=F(x^{\star})$ is finite.
\tag{A4}
\end{enumerate}

\section*{Main Convergence Theorem}
\begin{theorem}[Ergodic $O(1/k)$ rate]
\label{thm:rate}
Let $\{x_{k}\}$ be generated by \eqref{eq:pgm} under Assumptions (A1)–(A4) with a stepsize $\eta\in(0,1/L]$. Then for any $k\ge1$,
\[
F(x_{k})-F^{\star}\;\le\;\frac{\|x_{0}-x^{\star}\|^{2}}{2\eta\,k}.
\tag{T1}
\]
Consequently, $F(x_{k})\to F^{\star}$ as $k\to\infty$.
\end{theorem}

\section*{Theoretical Properties}
\begin{enumerate}[label=\arabic*.]
\item \textbf{Stability.} The iteration is a non‑expansive mapping when $\eta\le1/L$, guaranteeing that the distance to any fixed point never increases.
\item \textbf{Hyperparameter sensitivity.} The bound \eqref{T1} improves linearly with $\eta$; choosing the largest admissible $\eta=1/L$ yields the smallest constant.
\item \textbf{Comparison to baselines.} When $g\equiv0$, \eqref{eq:pgm} reduces to gradient descent and Theorem~\ref{thm:rate} recovers the classical $O(1/k)$ rate for smooth convex minimisation. When $g$ is the $\ell_{1}$‑norm, the method becomes the iterative soft‑thresholding algorithm (ISTA) with the same rate.
\end{enumerate}

\section*{Discussion}
The proof hinges on two elementary facts: (i) the descent lemma for $L$‑smooth functions and (ii) the \emph{firm non‑expansiveness} of the proximal operator. By coupling these inequalities we obtain a one‑step decrease that telescopes, leading directly to the $O(1/k)$ bound without any stochasticity or variance terms.

\newpage
\section*{Preliminaries (Definitions and Lemmas)}
\begin{definition}[Proximal operator]
For $\eta>0$ and a convex function $g$, the proximal mapping is
\[
\operatorname{prox}_{\eta g}(v)=\arg\min_{u}
\Big\{ g(u)+\frac{1}{2\eta}\|u-v\|^{2}\Big\}.
\]
\end{definition}

\begin{lemma}[Firm non‑expansiveness]
\label{lem:nonexp}
For any $v,w\in\mathbb{R}^{d}$,
\[
\|\operatorname{prox}_{\eta g}(v)-\operatorname{prox}_{\eta g}(w)\|^{2}
\le \langle \operatorname{prox}_{\eta g}(v)-\operatorname{prox}_{\eta g}(w),\,v-w\rangle .
\tag{L1}
\]
\end{lemma}
\begin{proof}
Standard property of proximal mappings; follows from monotonicity of the sub‑differential of $g$.
\end{proof}

\begin{lemma}[Descent lemma for $L$‑smooth $f$]
\label{lem:descent}
For all $x,y\in\mathbb{R}^{d}$,
\[
f(y)\le f(x)+\langle\nabla f(x),y-x\rangle+\frac{L}{2}\|y-x\|^{2}.
\tag{L2}
\]
\end{lemma}
\begin{proof}
Integrate the Lipschitz condition on $\nabla f$; see standard convex analysis references.
\end{proof}

\section*{Main Proof (Step‑by‑Step Derivation)}
Let $x^{\star}\in\arg\min F$ be an arbitrary optimal solution.  
Define the gradient step
\[
y_{k}=x_{k}-\eta\nabla f(x_{k}),\qquad x_{k+1}= \operatorname{prox}_{\eta g}(y_{k}).
\tag{P1}
\]

\begin{enumerate}[label=[Step \arabic*], leftmargin=*]
\item \textbf{Optimality condition of the proximal step.}  
By definition of $\operatorname{prox}_{\eta g}$,
\[
0\in \partial g(x_{k+1})+\frac{1}{\eta}(x_{k+1}-y_{k}),
\]
which is equivalent to
\[
\langle x_{k+1}-y_{k},\,u-x_{k+1}\rangle+\eta\bigl(g(x_{k+1})-g(u)\bigr)\ge0,
\quad\forall u\in\mathbb{R}^{d}.
\tag{S1}
\]

\item \textbf{Insert $y_{k}=x_{k}-\eta\nabla f(x_{k})$.}  
Replacing $y_{k}$ in \eqref{S1} yields
\[
\langle x_{k+1}-x_{k}+\eta\nabla f(x_{k}),\,u-x_{k+1}\rangle
+\eta\bigl(g(x_{k+1})-g(u)\bigr)\ge0.
\tag{S2}
\]

\item \textbf{Rearrange terms.}  
Expanding the inner product,
\[
\langle x_{k+1}-x_{k},\,u-x_{k+1}\rangle
+\eta\langle\nabla f(x_{k}),\,u-x_{k+1}\rangle
+\eta\bigl(g(x_{k+1})-g(u)\bigr)\ge0.
\tag{S3}
\]

\item \textbf{Apply the identity}
\[
\langle a-b,\,c-b\rangle
= \tfrac12\bigl(\|a-c\|^{2}-\|a-b\|^{2}-\|b-c\|^{2}\bigr),
\]
with $a=x_{k+1}$, $b=x_{k}$, $c=u$.  
Thus,
\[
\langle x_{k+1}-x_{k},\,u-x_{k+1}\rangle
= \tfrac12\Bigl(\|x_{k}-u\|^{2}-\|x_{k+1}-u\|^{2}-\|x_{k+1}-x_{k}\|^{2}\Bigr).
\tag{S4}
\]

\item \textbf{Insert \eqref{S4} into \eqref{S3}.}  
We obtain
\[
\frac12\Bigl(\|x_{k}-u\|^{2}-\|x_{k+1}-u\|^{2}-\|x_{k+1}-x_{k}\|^{2}\Bigr)
+\eta\langle\nabla f(x_{k}),\,u-x_{k+1}\rangle
+\eta\bigl(g(x_{k+1})-g(u)\bigr)\ge0.
\tag{S5}
\]

\item \textbf{Move the non‑negative term $ \|x_{k+1}-x_{k}\|^{2}$ to the right.}
\[
\frac12\bigl(\|x_{k}-u\|^{2}-\|x_{k+1}-u\|^{2}\bigr)
\le \frac12\|x_{k+1}-x_{k}\|^{2}
-\eta\langle\nabla f(x_{k}),\,u-x_{k+1}\rangle
-\eta\bigl(g(x_{k+1})-g(u)\bigr).
\tag{S6}
\]

\item \textbf{Bound the term $\langle\nabla f(x_{k}),\,u-x_{k+1}\rangle$.}  
Add and subtract $x_{k}$:
\[
\langle\nabla f(x_{k}),\,u-x_{k+1}\rangle
= \langle\nabla f(x_{k}),\,u-x_{k}\rangle
+ \langle\nabla f(x_{k}),\,x_{k}-x_{k+1}\rangle .
\tag{S7}
\]

\item \textbf{Apply Lemma~\ref{lem:descent} with $x=x_{k}$, $y=u$.}  
From \eqref{L2},
\[
f(u)\ge f(x_{k})+\langle\nabla f(x_{k}),\,u-x_{k}\rangle
-\frac{L}{2}\|u-x_{k}\|^{2}.
\tag{S8}
\]
Rearranging gives
\[
\langle\nabla f(x_{k}),\,u-x_{k}\rangle
\le f(u)-f(x_{k})+\frac{L}{2}\|u-x_{k}\|^{2}.
\tag{S9}
\]

\item \textbf{Bound $\langle\nabla f(x_{k}),\,x_{k}-x_{k+1}\rangle$ using Cauchy–Schwarz
and Young’s inequality.}  
For any $\alpha>0$,
\[
\langle\nabla f(x_{k}),\,x_{k}-x_{k+1}\rangle
\le \frac{1}{2\alpha}\|\nabla f(x_{k})\|^{2}
+\frac{\alpha}{2}\|x_{k}-x_{k+1}\|^{2}.
\tag{S10}
\]

\item \textbf{Choose $\alpha = L$ and combine \eqref{S7}–\eqref{S10}.}  
Using $\eta\le 1/L$, we have $\eta\alpha = \eta L\le1$. Hence,
\[
-\eta\langle\nabla f(x_{k}),\,u-x_{k+1}\rangle
\le -\eta\bigl[f(u)-f(x_{k})\bigr]
-\frac{\eta L}{2}\|u-x_{k}\|^{2}
+\frac{\eta}{2L}\|\nabla f(x_{k})\|^{2}
+\frac{\eta L}{2}\|x_{k+1}-x_{k}\|^{2}.
\tag{S11}
\]

\item \textbf{Insert \eqref{S11} into inequality \eqref{S6}.}  
After cancellation of the $\frac{\eta L}{2}\|x_{k+1}-x_{k}\|^{2}$ term with
the $\frac12\|x_{k+1}-x_{k}\|^{2}$ term (recall $\eta L\le1$), we obtain
\[
\frac12\bigl(\|x_{k}-u\|^{2}-\|x_{k+1}-u\|^{2}\bigr)
\le \eta\bigl[F(x_{k})-F(u)\bigr]
-\frac{\eta}{2L}\|\nabla f(x_{k})\|^{2}
+\frac12\|x_{k+1}-x_{k}\|^{2}.
\tag{S12}
\]

\item \textbf{Drop the non‑negative term $-\frac{\eta}{2L}\|\nabla f(x_{k})\|^{2}$.}  
Since it only strengthens the inequality,
\[
\frac12\bigl(\|x_{k}-u\|^{2}-\|x_{k+1}-u\|^{2}\bigr)
\le \eta\bigl[F(x_{k})-F(u)\bigr]
+\frac12\|x_{k+1}-x_{k}\|^{2}.
\tag{S13}
\]

\item \textbf{Bound $\|x_{k+1}-x_{k}\|^{2}$ using Lemma~\ref{lem:nonexp}.}  
From \eqref{L1} with $v=y_{k}$ and $w=x_{k}$,
\[
\|x_{k+1}-x_{k}\|^{2}
\le \langle x_{k+1}-x_{k},\,y_{k}-x_{k}\rangle
= -\eta\langle x_{k+1}-x_{k},\,\nabla f(x_{k})\rangle .
\tag{S14}
\]
Applying Young’s inequality with the same $\alpha=L$ yields
\[
\|x_{k+1}-x_{k}\|^{2}
\le \frac{\eta^{2}}{2L}\|\nabla f(x_{k})\|^{2}
+\frac{\eta L}{2}\|x_{k+1}-x_{k}\|^{2}.
\]
Since $\eta L\le1$, we can move the last term to the left and obtain
\[
\|x_{k+1}-x_{k}\|^{2}\le \frac{\eta^{2}}{L}\|\nabla f(x_{k})\|^{2}.
\tag{S15}
\]

\item \textbf{Insert \eqref{S15} into \eqref{S13}.}  
The term $\frac12\|x_{k+1}-x_{k}\|^{2}$ becomes at most
$\frac{\eta^{2}}{2L}\|\nabla f(x_{k})\|^{2}$, which we again drop (non‑negative). Hence
\[
\frac12\bigl(\|x_{k}-u\|^{2}-\|x_{k+1}-u\|^{2}\bigr)
\le \eta\bigl[F(x_{k})-F(u)\bigr].
\tag{S16}
\]

\item \textbf{Set $u=x^{\star}$ (an optimal solution).}  
Since $F(x^{\star})=F^{\star}$,
\[
\frac12\bigl(\|x_{k}-x^{\star}\|^{2}-\|x_{k+1}-x^{\star}\|^{2}\bigr)
\le \eta\bigl[F(x_{k})-F^{\star}\bigr].
\tag{S17}
\]

\item \textbf{Rearrange to isolate the function gap.}
\[
F(x_{k})-F^{\star}\ge \frac{1}{2\eta}
\bigl(\|x_{k}-x^{\star}\|^{2}-\|x_{k+1}-x^{\star}\|^{2}\bigr).
\tag{S18}
\]

\item \textbf{Sum \eqref{S18} from $k=0$ to $k=K-1$.}  
Telescoping yields
\[
\sum_{k=0}^{K-1}\bigl[F(x_{k})-F^{\star}\bigr]
\ge \frac{1}{2\eta}\bigl(\|x_{0}-x^{\star}\|^{2}-\|x_{K}-x^{\star}\|^{2}\bigr)
\ge \frac{1}{2\eta}\|x_{0}-x^{\star}\|^{2}.
\tag{S19}
\]

\item \textbf{Average the inequality.}  
Dividing by $K$ and using that $F(x_{k})\ge F^{\star}$,
\[
\frac{1}{K}\sum_{k=0}^{K-1}\bigl[F(x_{k})-F^{\star}\bigr]
\ge \frac{\|x_{0}-x^{\star}\|^{2}}{2\eta K}.
\tag{S20}
\]

\item \textbf{Since the sequence $\{F(x_{k})\}$ is non‑increasing,}
\[
F(x_{K})-F^{\star}\le \frac{1}{K}\sum_{k=0}^{K-1}\bigl[F(x_{k})-F^{\star}\bigr].
\tag{S21}
\]

\item \textbf{Combine \eqref{S20} and \eqref{S21}.}  
We finally obtain the desired bound
\[
F(x_{K})-F^{\star}\le \frac{\|x_{0}-x^{\star}\|^{2}}{2\eta K},
\qquad\forall\,K\ge1,
\]
which is exactly \eqref{T1}.
\end{enumerate}

\section*{Rate Derivation (With Constants)}
The constant appearing in \eqref{T1} is $C:=\frac{\|x_{0}-x^{\star}\|^{2}}{2\eta}$.
With the maximal admissible stepsize $\eta=1/L$, we have
\[
F(x_{k})-F^{\star}\le \frac{L\|x_{0}-x^{\star}\|^{2}}{2k},
\]
showing explicitly the dependence on the smoothness parameter $L$.

\section*{Special Cases}
\begin{itemize}
\item \textbf{Pure gradient descent.} $g\equiv0$ $\Rightarrow$ $\operatorname{prox}_{\eta g}=\mathrm{Id}$, and \eqref{T1} reduces to the classical $O(1/k)$ rate for smooth convex minimisation.
\item \textbf{Lasso (elastic net) regularisation.} $g(x)=\lambda\|x\|_{1}$, the proximal step is soft‑thresholding. The same bound \eqref{T1} holds, providing convergence guarantees for ISTA.
\item \textbf{Indicator function of a convex set $C$.} $g=\iota_{C}$, then $\operatorname{prox}_{\eta g}$ is the Euclidean projection onto $C$. The method becomes projected gradient descent with the same $O(1/k)$ guarantee.
\end{itemize}

\end{document}